\section{问题分析}

\subsection{字段语义与数据口径}

住院、门诊、体检三张明细均以项目行为基本单位，同一患者同一开单时刻可以出现多行。若直接把明细行视为患者，会重复计算多项目申请；若只按患者编号聚合，又会把同一患者不同时期的多次申请合并。因此定义申请事件键
\begin{equation}
  e=(s,p,a),
\end{equation}
其中$s$为来源，$p$为患者编号，$a$为精确开单时刻；患者编号缺失时以来源行号构造只在该表内唯一的代理键。对同一事件内完全相同的服务项目去重，而不同项目全部保留。打印、预约和图文报告等不占用超声扫查设备的记录被标记为非服务项，不进入工作量。

时间字段分为到达与完成两类。住院释放时刻取开单时刻$a_e$，截止时刻为$d_e=a_e+48\unit{h}$。门诊和体检缺少真实预约日期，本文只在回顾性评价中使用报告日期构造计划日$g_e$；其日内窗口为上午$[8{:}00,12{:}00]$和下午$[13{:}00,17{:}00]$。历史报告的具体时分、历史机房和历史医生不进入决策。官方质量指标强调以申请单为分母评价完成情况\cite{nhc2022}，但附件又缺少独立的报告提交时刻，因此本文将排程中最后一个项目结束时刻作为“完成检查并可进入报告环节”的代理，全文不把它表述为真实报告提交时间。

\begin{table}[H]
  \centering
  \caption{统一连续模拟与留出期评价规模}
  \label{tab:data-scope}
  \begin{tabular}{lrrr}
\toprule
患者来源 & 连续模拟事件数 & 项目任务数 & 留出期评价事件数 \\
\midrule
住院 & 71,029 & 141,429 & 70,771 \\
门诊 & 55,512 & 69,756 & 55,226 \\
体检 & 40,820 & 117,701 & 40,820 \\
\midrule
合计 & 167,361 & 328,886 & 166,817 \\
\bottomrule
\end{tabular}

\end{table}

连续模拟规模见表\ref{tab:data-scope}。其中住院事件在2日预热期和一年留出期均保留精确开单时刻；门诊与体检只使用计划日压力，不使用历史派单结果。这一处理使问题三能够使用三类患者的真实项目组合和日需求，同时避免把不可获得的未来机房、医生当作模型输入。

\subsection{项目归类与资源标记}

项目归类采用互斥优先级词典。首先识别非服务项；随后依次识别产科III/IV级、产科I/II级、心脏、介入/定位、血管、腔内、腹部/胃肠、浅表器官、儿科专项和一般其他。优先级解决关键词重叠，例如“心脏造影”先被识别为心脏项目，不因包含“造影”而误入介入类别。每个项目族的时长边界直接依据题给区间，并保留项目原名用于设备匹配，不以类别名替代项目名。

题目中的检查准备要求转换为三个事件级标记：包含“空腹”的项目必须在10时前完成；包含“充盈膀胱”等词的事件在主情景中以开单后60分钟作为准备完成时刻，另以45和90分钟作敏感性；若门诊/体检在计划日开班前已满足提前量，则从工作块起点即可排入；明确“床旁”的检查只允许7号床边B超机。若一个事件同时包含多个标记，所有相应约束共同生效。

\subsection{时间边界和中文名称泄漏控制}

用于建模的项目目录只统计训练期出现的项目名称。留出期项目在调度输入中按训练期确定的归类规则处理；若新名称无法获得项目级设备证据，才进入已标注的类别回退。另输出全期项目目录只用于检查名称漂移，不参与留出期派位。医生效率和报告间隔同样只使用2024年3月31日以前的报告记录，避免利用留出期完成时刻修正未来调度速度。

\subsection{需求变化与高峰特征}

图\ref{fig:monthly-demand}给出三类患者月度日均开单量。住院需求总体平稳上升；门诊呈较强年度和节假日波动；体检最不稳定，在团检月份形成尖峰。因而用全年平均值配置日容量会同时低估体检高峰和住院周一压力。本文在问题一中分别计算来源内日中位数、90\%分位和95\%分位，并保留星期结构，不把“高峰”定义为简单事件数最多的若干天。

\begin{figure}[H]
  \centering
  \includegraphics[width=\textwidth]{fig01_monthly_demand.pdf}
  \caption{三类患者月度日均开单量及时间顺序留出边界}
  \label{fig:monthly-demand}
\end{figure}

训练期科室需求见表\ref{tab:department}. 住院需求分散于87个可识别科室，前6科室占35.45\%；门诊需求集中度更高，前6科室占62.95\%；体检记录只有一个科室编码，不能据此推断其临床科室组成。科室字段缺失的事件数量很少，但模型仍在结果中保留“缺失”类别，不使用众数强行填补。

\begin{table}[H]
  \centering
  \caption{训练期科室需求集中度}
  \label{tab:department}
  \begin{tabular}{lrrrr}
\toprule
来源 & 训练期事件数 & 可识别科室数 & 前6科室占比 & 前12科室占比 \\
\midrule
住院 & 297,656 & 87 & 35.45\% & 56.72\% \\
门诊 & 245,614 & 133 & 62.95\% & 81.33\% \\
体检 & 239,569 & 1 & 100.00\% & 100.00\% \\
\bottomrule
\end{tabular}

\end{table}

\subsection{数据能够支持与不能支持的结论}

数据能够支持：事件到达规模、项目组合、题给时长边界、项目文本对应的设备能力、历史医生活动时段以及留出期回顾性压力测试。数据不能直接支持：真实扫查时长、报告撰写时长、未来医生身份与项目资质、真实预约时刻、未完成患者的真实后续处置。过程挖掘中的相邻时间戳不自动等于活动持续时间\cite{vda2011}，影像报告时间的测量也必须区分图像采集、阅片和报告环节\cite{sexauer2022}。因此，本文所有“时长估计”均称为调度占用参数，所有“完成率”均注明代理边界。
